% Einheit 5 von 5 – Zusammengesetzte Körper und Anwendungen (bank/koerper/e5.jsonl, zusammenbau v0.1)
%% TODO Zweigzeile Teil 1: was hier gelernt wird (Satz in Schülersprache aus dem Einheitstitel) – Einheit 5
\einheitenkopf[e5]{Einheit 5 von 5 · Zusammengesetzte Körper und Anwendungen}
\zweigzeile{ab Kl. 6, je nach Buch bis Kl. 9 (am Gymnasium ab Kl. 9) · P10}
\setcounter{aufgabe}{24}

%% TODO Titel: Ich-kann-Satz fehlt in der Bank (Platzhalter: Kettenname) – Nr. 25
%% TODO Anweisung: ein Satz über der ersten Teilaufgabe fehlt in der Bank – Nr. 25
\begin{aufgabe}{Zusammengesetzt}
\begin{teile}
\teil Ein Haus mit Satteldach: Aus welchen zwei Körpern besteht es?
\end{teile}
%% TODO Darstellung für schwach fehlt in der Bank (koerper-e5-k1-s1-v1; Abschnitt „Für schwache Schüler“)
\swz{Eine Betonstufe besteht aus zwei Quadern: unten $8$ dm × $4$ dm × $2$ dm, oben darauf $8$ dm × $2$ dm × $2$ dm. Volumen?}{}
%% TODO Darstellung für schwach fehlt in der Bank (koerper-e5-k1-s1-v2; Abschnitt „Für schwache Schüler“)
\swz{Ein Regal in L-Form besteht aus zwei Holzblöcken: $9$ cm × $3$ cm × $2$ cm und $3$ cm × $3$ cm × $5$ cm. Volumen?}{}
%% TODO Darstellung für schwach fehlt in der Bank (koerper-e5-k1-s1-v3; Abschnitt „Für schwache Schüler“)
\swz{Ein Büro aus zwei Containern: unten $7$ m × $3$ m × $3$ m, oben $4$ m × $3$ m × $2$ m. Volumen?}{}
%% TODO Darstellung für schwach fehlt in der Bank (koerper-e5-k1-s1-v4; Abschnitt „Für schwache Schüler“)
\swz{Ein Treppenpodest aus zwei Quadern: $5$ m × $2$ m × $1$ m und $5$ m × $1$ m × $2$ m. Volumen?}{}
%% TODO Darstellung für schwach fehlt in der Bank (koerper-e5-k1-s1-v5; Abschnitt „Für schwache Schüler“)
\swz{Ein Sockel aus zwei Quadern: unten $9$ dm × $6$ dm × $1$ dm, darauf $3$ dm × $3$ dm × $7$ dm. Volumen?}{}
\swfrage{Was bleibt gleich, was ändert sich?}
%% TODO Darstellung für schwach fehlt in der Bank (koerper-e5-k1-s2-v1; Abschnitt „Für schwache Schüler“)
\swz{Ein Haus: Quader $11$ m lang, $7$ m breit, $3$ m hoch; darauf ein Satteldach, $7$ m breit, $3$ m hoch, so lang wie das Haus – Volumen?}{}
%% TODO Darstellung für schwach fehlt in der Bank (koerper-e5-k1-s3-v1; Abschnitt „Für schwache Schüler“)
\swz{Ein Pool: Rechteck $8$ m × $4$ m, an der $4$-m-Seite ein Halbkreis mit $2$ m Radius; überall $1{,}4$ m tief – Wasservolumen auf eine Stelle?}{}
\begin{teile}
\teil Ein Pool: Rechteck $7$ m × $4$ m, an der $4$-m-Seite ein Halbkreis mit $2$ m Radius, überall $1{,}6$ m tief. Welcher Term gibt das Wasservolumen in m³? \\ \kreuz{$7 \cdot 4 + \tfrac{1}{2} \cdot \pi \cdot 2^2 \cdot 1{,}6$} \\ \kreuz{$(7 \cdot 4 + \tfrac{1}{2} \cdot \pi \cdot 2^2) \cdot 1{,}6$} \\ \kreuz{$(7 \cdot 4 + \pi \cdot 2^2) \cdot 1{,}6$}
\end{teile}
%% TODO Darstellung für schwach fehlt in der Bank (koerper-e5-k1-s5-v1; Abschnitt „Für schwache Schüler“)
\swz{Ein Tortenkarton ist innen $24$ cm × $24$ cm × $9$ cm groß. Darin steht eine runde Torte mit $22$ cm Durchmesser und $8$ cm Höhe. Wie viel Luft bleibt im Karton, auf eine Stelle?}{}
\begin{teile}
\teil Ein zylinderförmiger Pflanzkübel (Radius $18$ cm, Höhe $55$ cm) soll in einem quaderförmigen Karton verschickt werden, der so klein wie möglich ist. Kreuze an und begründe \hfill (P10 2024 OS). \\ \kreuz{$19$ cm × $19$ cm × $56$ cm} \\ \kreuz{$37$ cm × $37$ cm × $56$ cm} \\ \kreuz{$37$ cm × $37$ cm × $53$ cm} \\ \kreuz{$56$ cm × $56$ cm × $56$ cm}
\end{teile}
%% TODO Darstellung für schwach fehlt in der Bank (koerper-e5-k1-s7-v1; Abschnitt „Für schwache Schüler“)
\swz{Eine Kiste ist innen $50$ cm × $34$ cm × $21$ cm groß. Hinein sollen Schachteln mit $11$ cm × $8$ cm × $7$ cm, alle gleich ausgerichtet ($11$ cm entlang der $50$ cm, $8$ cm entlang der $34$ cm). Wie viele passen hinein?}{}
%% TODO Darstellung für schwach fehlt in der Bank (koerper-e5-k1-s8-v1; Abschnitt „Für schwache Schüler“)
\swz{Ein Aquarium ist innen $90$ cm × $35$ cm × $42$ cm groß und randvoll. $1$ dm³ Wasser wiegt $1$ kg. Wie schwer ist das Wasser?}{}
\swz[3]{Ein kegelförmiger Lampenschirm (Radius $18$ cm, Höhe $35$ cm) liegt in einer zylinderförmigen Verpackung (innen Radius $18{,}5$ cm, Höhe $36$ cm). Für den Kegel gilt $V = \tfrac{1}{3} \cdot \pi \cdot r^2 \cdot h$ ($r$ Radius, $h$ Höhe; Formelsammlung). Berechne das Restvolumen der Verpackung in cm³ \hfill (P10 2024 OS).}{}
\end{aufgabe}

%% TODO Titel: Ich-kann-Satz fehlt in der Bank (Platzhalter: Kettenname) – Nr. 26
%% TODO Anweisung: ein Satz über der ersten Teilaufgabe fehlt in der Bank – Nr. 26
\begin{aufgabe}{Füllstand in Prozent}
%% TODO Darstellung für schwach fehlt in der Bank (koerper-e5-k2-s1-v1; Abschnitt „Für schwache Schüler“)
\swz{Ein Aquarium fasst $250$ l; es sind $185$ l Wasser darin. Füllstand in Prozent?}{}
\end{aufgabe}

%% TODO Titel: Ich-kann-Satz fehlt in der Bank (Platzhalter: Kettenname) – Nr. 27
%% TODO Anweisung: ein Satz über der ersten Teilaufgabe fehlt in der Bank – Nr. 27
\begin{aufgabe}{Oberfläche eines zusammengesetzten Körpers}
%% TODO Darstellung für schwach fehlt in der Bank (koerper-e5-k3-s1-v1; Abschnitt „Für schwache Schüler“)
\swz{Zwei Würfel mit $5$ cm Kantenlänge werden aufeinandergeklebt. Wie groß ist die Oberfläche des neuen Körpers?}{}
\end{aufgabe}

%% TODO Titel: Ich-kann-Satz fehlt in der Bank (Platzhalter: Kettenname) – Nr. 28
%% TODO Anweisung: ein Satz über der ersten Teilaufgabe fehlt in der Bank – Nr. 28
\begin{aufgabe}{Zusammengesetzt – Fehler finden, Begründen, Anwendung, Darstellungswechsel}
\swz[3]{Ein Becken: Rechteck $5$ m × $4$ m, an der $4$-m-Seite ein Halbkreis mit $2$ m Radius, überall $1{,}7$ m tief. Ben rechnet: \rechnung{5 \cdot 4 + \tfrac{1}{2} \cdot \pi \cdot 2^2 \cdot 1{,}7 &\approx 30{,}7 \text{ m}^3} Wo ist der Fehler?}{}
\swfrage{Erkläre, ohne zu rechnen, wie man das Volumen eines Hauses mit Satteldach bestimmt.}
\swz[3]{Ein Pool: Rechteck $9$ m × $4$ m, an der $4$-m-Seite ein Halbkreis mit $2$ m Radius, überall $1{,}4$ m tief. Der Schlauch liefert $1\,800$ l pro Stunde. Ist der Pool am Sonntag um $14$ Uhr voll, wenn man am Freitag um $18$ Uhr anfängt?}{}
\swz{Ein Haus: Quader $7$ m lang, $5$ m breit, $3$ m hoch; darauf ein Satteldach, $5$ m breit, $2$ m hoch, $7$ m lang. Stelle einen Term für das Volumen auf und berechne ihn.}{}
\end{aufgabe}

\uebersichtskasten{Zerlegen: Körper in Quader, Prismen, Zylinder oder halbe Zylinder schneiden, jedes Volumen berechnen, addieren – oder vom großen Körper den fehlenden Teil abziehen. \\ Haus: Quader 8 m · 6 m · 3 m und Dach als Dreiecksprisma (6 m breit, 2 m hoch, 8 m lang): 144 + 48 = 192 m³ \\ Term prüfen: Bei „(Rechteck + Halbkreis) mal Tiefe“ muss die Klammer stehen – sonst wirkt die Tiefe nur auf einen Teil. \\ Passt es hinein? Der Durchmesser (nicht der Radius) muss in die Breite passen; beim Stapeln anordnen, nicht nur Volumen teilen. \\ Masse aus Volumen: Volumen in dm³ mal Dichte (Wasser: 1 kg je dm³).}
